%! Graphing Rational Functions — Unit 5, Lesson 5.5 — Warm-Up (Answer Key)
\documentclass[10pt]{article}
\usepackage{algebra2-article}
\usepackage{algebra2-key}   % pulls in -boxes; do NOT also load -boxes
\newcommand{\TallMath}[1]{$\displaystyle #1\rule[-1.4em]{0pt}{3.2em}$}

\begin{document}

\pageheader{Unit 5, Lesson 5.5}{Warm-Up --- Answer Key}
\namedateperiod

\vspace{0.10in}

\begin{notesbox}{Getting Ready: Skills We'll Use Today}
\small Three skills. Today's whole lesson is these three, run in order, on one function at a time.

\vspace{4pt}
\textbf{1. Factor, restrict, cancel, sort} (Lessons 5.0 and 5.1). \; $h(x)=\dfrac{x^2-9}{x^2-2x-3}$
\begin{enumerate}[label=(\alph*), leftmargin=2.1em, itemsep=3pt]
  \item Factor both: \; $h(x)=\dfrac{(\text{\ans{$x-3$}})(\text{\ans{$x+3$}})}%
        {(\text{\ans{$x-3$}})(\text{\ans{$x+1$}})}$
  \item Restrictions, read off the \emph{original} denominator: $x\neq$ \ans{$3$} and
        $x\neq$ \ans{$-1$}.
  \item Cancel the common factor. Simplified: \ans{$\frac{x+3}{x+1}$}
  \item One restriction cancelled and one did not. The one that cancelled gives a
        \ans{hole} at $x=$ \ans{$3$}; the one that survived gives a \ans{vertical asymptote} at
        $x=$ \ans{$-1$}.
\end{enumerate}

\vspace{4pt}
\textbf{2. The degree comparison} (Lesson 5.0). Give the horizontal asymptote --- or say there is none.
\begin{enumerate}[label=(\alph*), leftmargin=2.1em, itemsep=3pt]
  \item $y=\dfrac{2x}{x^2+5}$: \; top degree \ans{$1$}, bottom degree \ans{$2$}, so
        HA \ans{$y=0$}
  \item $y=\dfrac{3x^2-1}{x^2+4}$: \; top degree \ans{$2$}, bottom degree \ans{$2$}, so
        HA \ans{$y=3$}
  \item $y=\dfrac{x^3}{x^2-1}$: \; top degree \ans{$3$}, bottom degree \ans{$2$}, so
        HA \ans{none}
\end{enumerate}

\vspace{4pt}
\textbf{3. Signs, not values} (Lesson 4.6's sign chart). For $f(x)=\dfrac{x+1}{(x-2)(x+3)}$, decide
whether each output is \emph{positive} or \emph{negative}. Do not compute the value --- just track the
sign of each factor.
\begin{center}
\renewcommand{\arraystretch}{1.5}
\begin{tabularx}{0.94\linewidth}{|C{1.6cm}|C{2.0cm}|C{2.0cm}|C{2.0cm}|X|}
\hline
\rowcolor{forestbg}
\textbf{$x$} & \textbf{sign of $x+1$} & \textbf{sign of $x-2$} & \textbf{sign of $x+3$} &
\textbf{sign of $f(x)$} \\ \hline
$-4$ & \ans{$-$} & \ans{$-$} & \ans{$-$} & \ans{negative} \\ \hline
$-2$ & \ans{$-$} & \ans{$-$} & \ans{$+$} & \ans{positive} \\ \hline
$0$  & \ans{$+$} & \ans{$-$} & \ans{$+$} & \ans{negative} \\ \hline
$3$  & \ans{$+$} & \ans{$+$} & \ans{$+$} & \ans{positive} \\ \hline
\end{tabularx}
\end{center}
\vspace{-2pt}
You never needed the actual outputs --- only the \ans{sign} of each factor.

\end{notesbox}

\begin{teachernote}
One item per leg of the lesson, and today the three legs are literally the three steps of the procedure.

\vspace{3pt}
\textbf{Item 1} is the cancel test from 5.0, driven by 5.1's algebra. The line to insist on is 1(b): the
restrictions come off the \emph{original} denominator, before any cancelling. A student who simplifies
first and then restricts loses the hole entirely --- and that same error becomes an extraneous solution
in 5.6.

\vspace{3pt}
\textbf{Item 2} is pure recall, and 2(c) is the one that matters today: when the top outranks the bottom
there is \emph{no} horizontal asymptote, and ``none'' is a complete answer. (A student who asks what the
graph does instead has asked the Tier E question --- the answer is a slant asymptote, which is
enrichment only and is never assessed.)

\vspace{3pt}
\textbf{Item 3} is the new skill, disguised as review. Students have built sign charts since 4.6, but
always for polynomials; the only change today is that one of the factors sits in the denominator, and a
quotient of signs behaves exactly like a product of signs. Debrief the last line out loud --- the whole
point of a sign chart is that it is cheaper than evaluating. Note the alternation down the last column
($-,+,-,+$): every boundary here is a single factor, so the sign flips at each one. Do not promise that
it always alternates; that promise breaks the first time a factor is squared.
\end{teachernote}

\end{document}
